% Supplementary material for the AJP teaching article in main.tex. Build main.tex first: the references to its
% equations and sections are read from main.aux through xr-hyper.
% Build with:  latexmk -pdf main.tex      (from this folder)
% Figures:     Rscript paper/figures/make_figures.R      (from the repository root)
% To estimate the length in AJP's two-column journal format, replace "preprint" by "reprint".
\documentclass[aps,prb,preprint,amsmath,amssymb,nofootinbib,longbibliography]{revtex4-2}


\usepackage{graphicx}
\usepackage{bm}
\usepackage{braket}
\usepackage{xcolor}
\usepackage{algpseudocode}
\usepackage{xr-hyper}
\usepackage[colorlinks=true,allcolors=blue!50!black]{hyperref}

\newcommand{\todo}[1]{\textcolor{red}{[TODO: #1]}}
\newcommand{\ev}[1]{\left\langle #1 \right\rangle}
\newcommand{\Eloc}{E_{\mathrm{loc}}}
\newcommand{\dtau}{\delta\tau}
\renewcommand{\Re}{\operatorname{Re}}
\renewcommand{\Im}{\operatorname{Im}}
\externaldocument[M-]{main}
% Shared by main.tex and supplementary.tex. AJP reviews anonymously: the first submission leaves out the author,
% affiliation and acknowledgments, and the citation of the code, which names the authors. Set \anonymousfalse for
% the version with them (after conditional acceptance).
\newif\ifanonymous
\anonymoustrue
\newcommand{\codecite}{\ifanonymous\cite{nnqs_code_anon}\else\cite{nnqs_code}\fi}

\renewcommand{\thefigure}{S\arabic{figure}}
\renewcommand{\thetable}{S\arabic{table}}
\renewcommand{\theequation}{S\arabic{equation}}
\renewcommand{\thesection}{S\arabic{section}}

\begin{document}

\title{Supplementary material for ``Neural-network quantum states as imaginary-time evolution: a worked example on the transverse-field Ising chain''}
\date{\today}
\maketitle

\section{Implementation details}
\label{sec:impl}

Figure~\ref{alg:training} lists one training run of the program\codecite\ described in Sec.~\ref{M-sec:implementation} of the
main text, and Table~\ref{tab:settings} gives its settings. The comparisons of the sign gauge (Sec.~\ref{M-sec:model}),
of the full and block-diagonal $S$ (Sec.~\ref{M-sec:critical}) and of the implementation errors
(Sec.~\ref{M-sec:pitfalls}) use a single image; all other results use four.

\paragraph*{Sampling.}
Each image runs one Markov chain for the whole training, continuing from where the previous epoch left off. After
each parameter update it takes $5n$ steps to re-equilibrate, and then records a sample after every sweep of $n$
proposals. The vector $\theta(s)$ is cached and updated in $\mathcal{O}(m)$ operations per accepted flip, but it is
recomputed from scratch after each parameter update.

\paragraph*{The SR step.}
For each block, the program centers the sampled $O_k$ and $\Eloc$. It then forms the block of $S$ and $F$ with the
unbiased normalization $1/(N_s-1)$, adds $\epsilon$ to the diagonal, and solves the Hermitian positive-definite
system with LAPACK's packed-storage driver \texttt{ppsvx}. In the code the hidden biases are stored as $b = c^*$, so
their update is conjugated.

\begin{figure}[h]
\begin{algorithmic}[1]
\State $a, c \gets 0$;\; $W_{ij} \gets$ small, nearly uniform (Table~\ref{tab:settings});\; $s \gets$ random
\For{epoch $= 1, \ldots, 1000$}
  \State $\alpha \gets$ average of $\alpha$ over all images
  \State $\theta \gets c + W s$ \Comment{recompute after every update}
  \State $5n$ Metropolis steps \Comment{re-equilibrate}
  \For{$i = 1,\ldots,N_s$}
    \State $n$ Metropolis steps; \; $s_i \gets s$; \; compute $\Eloc(s_i)$
  \EndFor
  \State $\bar E \gets$ mean of $\Re\Eloc$, averaged over images
  \State form $O_k(s_i)$, then $S$ and $F$ for each block
  \State $\alpha_B \gets \alpha_B - \dtau\,(S_{BB} + \epsilon)^{-1} F_B$ for each block $B$
\EndFor
\end{algorithmic}
\caption{One training run, as carried out on each of the $P$ parallel images.}
\label{alg:training}
\end{figure}

\begin{table}[h]
\caption{Settings for all runs. The comparisons in Secs.~\ref{M-sec:model}, \ref{M-sec:critical} and
\ref{M-sec:pitfalls} of the main text use one image instead of four.}
\label{tab:settings}
\begin{ruledtabular}
\begin{tabular}{ll}
Samples per image, $N_s$ & 256 \\
Images, $P$ & 4 \\
Epochs & 1000 \\
Time step, $\dtau$ & 0.01 \\
Diagonal shift, $\epsilon$ & $10^{-3}$ \\
Burn-in after each update & $5n$ steps \\
Steps between samples & $n$ (one sweep) \\
Initial $a$, $c$ & 0 \\
Initial $\Re W_{ij}$ & normal, mean $0.01$, s.d.\ $10^{-4}$ \\
Initial $\Im W_{ij}$ & normal, mean $-0.005$, s.d.\ $10^{-5}$ \\
Floating point & single precision \\
\end{tabular}
\end{ruledtabular}
\end{table}

\section{Exact solution of the open chain}
\label{sec:exact}

We outline the derivation; Refs.~\onlinecite{lieb1961,pfeuty1970,mbeng2024} give the details. By the sign gauge of
Sec.~\ref{M-sec:model} of the main text we may take $J, B > 0$ and restore the stagger of
Eq.~(\ref{M-eq:stagger}) at the end.

\paragraph*{Fermions.}
Exchanging the roles of $\sigma^x$ and $\sigma^z$ (a rotation of the spin axes) gives $H = -J\sum_{j=1}^{n-1}
\sigma^x_j\sigma^x_{j+1} - B\sum_{j=1}^n \sigma^z_j$, and turns the correlation $\ev{\sigma^z_i\sigma^z_k}$ of the
main text into $\ev{\sigma^x_i\sigma^x_k}$. The Jordan--Wigner transformation\cite{jordan1928}
\begin{equation}
  \sigma^z_j = 1 - 2c_j^\dagger c_j, \qquad \sigma^x_j = \prod_{l<j}(1 - 2c_l^\dagger c_l)\,(c_j + c_j^\dagger)
\end{equation}
maps the spins onto fermions $c_j$. In terms of the $2n$ Majorana operators
\begin{equation}
  a_j = c_j + c_j^\dagger, \qquad b_j = i(c_j^\dagger - c_j), \qquad \{\gamma_p, \gamma_q\} = 2\delta_{pq},
\end{equation}
ordered as $\gamma = (a_1, b_1, a_2, b_2, \ldots, a_n, b_n)$, the strings cancel between neighbors and
\begin{equation}
  \sigma^z_j = -i\,a_j b_j, \qquad \sigma^x_j\sigma^x_{j+1} = -i\,b_j a_{j+1} .
  \label{eq:majorana-spins}
\end{equation}
The Hamiltonian is therefore quadratic,
\begin{equation}
  H = iB\sum_{j=1}^n a_j b_j + iJ\sum_{j=1}^{n-1} b_j a_{j+1} = \frac{i}{4}\sum_{p,q=1}^{2n} A_{pq}\gamma_p\gamma_q ,
\end{equation}
with the real antisymmetric $2n\times 2n$ matrix whose only nonzero entries are
\begin{equation}
  A_{a_j b_j} = -A_{b_j a_j} = 2B, \qquad A_{b_j a_{j+1}} = -A_{a_{j+1} b_j} = 2J ,
  \label{eq:majorana-matrix}
\end{equation}
where $A_{a_j b_j}$ means the entry in the rows and columns of $a_j$ and $b_j$.

\paragraph*{Energy.}
The matrix $iA$ is Hermitian. Because $A$ couples each $a$ only to $b$'s, its eigenvalues come in pairs $\pm
2\lambda_k$, where the $\lambda_k$ are the singular values of the $a$--$b$ block of $A/2$, which has $B$ on its
diagonal and $-J$ next to it. Signs and transposition do not change singular values, so the $\lambda_k$ are those of
the matrix $M$ of the main text. A Bogoliubov transformation to the eigenvectors of $iA$ gives $H = \sum_k
2\lambda_k\,(f_k^\dagger f_k - \tfrac12)$ with independent fermion modes $f_k$. The ground state empties every mode,
so $E_0 = -\sum_k\lambda_k$, which is Eq.~(\ref{M-eq:exact-energy}) of the main text.

\paragraph*{Correlations.}
In the ground state,
\begin{equation}
  \ev{\gamma_p\gamma_q} = \delta_{pq} + G_{pq}, \qquad G = \operatorname{sign}(iA) = U\operatorname{sign}(\Lambda)U^\dagger,
  \label{eq:majorana-G}
\end{equation}
where $iA = U\Lambda U^\dagger$ is the eigendecomposition and the sign is taken eigenvalue by eigenvalue. $G$ is purely
imaginary, and since the sign is an odd function, $G$ inherits the structure of $A$: it couples $a$'s only to $b$'s,
so $\ev{a_pa_q} = \ev{b_pb_q} = \delta_{pq}$. For $i < k$, Eq.~\eqref{eq:majorana-spins} gives
\begin{equation}
  \sigma^x_i\sigma^x_k = \prod_{j=i}^{k-1}\sigma^x_j\sigma^x_{j+1} = \prod_{j=i}^{k-1}\left(-i\,b_j a_{j+1}\right),
\end{equation}
a product of $2(k-i)$ Majorana operators. Wick's theorem writes its average as a Pfaffian of contractions. Only the
pairings of a $b$ with an $a$ survive, and the Pfaffian reduces to a determinant:
\begin{equation}
  \ev{\sigma^z_i\sigma^z_k} = \det\big[C_{pq}\big]_{p,q = i}^{k-1}, \qquad C_{pq} = \ev{-i\,b_p a_{q+1}}
  = -i\,G_{b_p a_{q+1}} .
  \label{eq:majorana-det}
\end{equation}
The matrix $C$ is real. For $J < 0$ the result is multiplied by $(-1)^{k-i}$.

\paragraph*{Numerics.}
One Hermitian eigendecomposition of the $2n\times 2n$ matrix $iA$ gives $G$, and each correlation is then a
$(k-i)\times(k-i)$ determinant, which is easily affordable for $n = 1024$. The function \texttt{exact\_correlations}
in \texttt{data\_visualization.R}, which comes with the program,\codecite\ implements Eqs.~\eqref{eq:majorana-matrix}--\eqref{eq:majorana-det} line by line.
For 8 spins it agrees with exact diagonalization to $10^{-15}$.

\bibliographystyle{ajp}
\bibliography{references}

\end{document}
